\documentclass[11pt]{article}
\usepackage[margin=2.5cm]{geometry}
\usepackage{amsmath, amssymb}
\usepackage{enumitem}
\usepackage{xcolor}
\usepackage{hyperref}

\title{Homework Solutions: Lectures 5 \& 6\\[4pt]
\large MLE \& MAP Estimation}
\author{Probability \& Statistics Course}
\date{}

\begin{document}
\maketitle

% ============================================================
\section*{Estimating a Sensor's True Reading (Normal--Normal)}
% ============================================================

\subsection*{Setup}

A thermometer measures temperature with known measurement noise. Each reading is:
\[
  X_i \mid \mu \sim N(\mu,\; \sigma^2), \qquad \sigma = 2\;{}^\circ\text{C} \;\text{(known)}
\]

You take $n = 5$ independent readings:
\[
  x_1 = 21.3, \quad x_2 = 19.8, \quad x_3 = 22.1, \quad x_4 = 20.5, \quad x_5 = 23.0
\]

The manufacturer says these sensors are calibrated around $20\;{}^\circ\text{C}$, so you use a Gaussian prior:
\[
  \mu \sim N(m,\; \tau^2) = N(20,\; 3^2)
\]

\textbf{Goal:} Estimate the true temperature $\mu$ using (a) MLE and (b) MAP.

% ----------------------------------------
\subsection*{Part (a): Maximum Likelihood Estimate}

\subsubsection*{Step 1: Write the likelihood}

Since the observations are independent and $\sigma^2$ is known, the likelihood depends only on $\mu$:
\[
  L(\mu) = \prod_{i=1}^n f(x_i \mid \mu) = \prod_{i=1}^n \frac{1}{\sqrt{2\pi}\,\sigma} \exp\!\left(-\frac{(x_i - \mu)^2}{2\sigma^2}\right)
\]

\subsubsection*{Step 2: Log-likelihood}
\[
  \ell(\mu) = \ln L(\mu) = -\frac{n}{2}\ln(2\pi\sigma^2) - \frac{1}{2\sigma^2}\sum_{i=1}^n (x_i - \mu)^2
\]

Only the second term depends on $\mu$. Maximizing $\ell(\mu)$ is equivalent to minimizing $\sum(x_i - \mu)^2$.

\subsubsection*{Step 3: Differentiate and solve}
\[
  \frac{d\ell}{d\mu} = \frac{1}{\sigma^2}\sum_{i=1}^n (x_i - \mu) = \frac{1}{\sigma^2}\left(\sum x_i - n\mu\right) \stackrel{!}{=} 0
\]

\[
  \boxed{\hat\mu_{\text{MLE}} = \bar{x} = \frac{1}{n}\sum_{i=1}^n x_i}
\]

\subsubsection*{Step 4: Compute}
\[
  \hat\mu_{\text{MLE}} = \frac{21.3 + 19.8 + 22.1 + 20.5 + 23.0}{5} = \frac{106.7}{5} = 21.34\;{}^\circ\text{C}
\]

\textbf{Note:} The MLE uses \textit{only} the data. It completely ignores the manufacturer's calibration information that the sensor should read around $20\;{}^\circ\text{C}$.


% ----------------------------------------
\subsection*{Part (b): MAP Estimate}

\subsubsection*{Step 1: Bayes' theorem --- what it says and why}

Bayes' theorem tells us how to \textit{update} our belief about $\mu$ after seeing data. Start from the definition of conditional probability:
\[
  p(\mu \mid \mathbf{x}) = \frac{p(\mathbf{x} \mid \mu)\; p(\mu)}{p(\mathbf{x})}
\]

Let's name each piece:

\begin{center}
\renewcommand{\arraystretch}{1.5}
\begin{tabular}{lcl}
$p(\mu)$ & = & \textbf{Prior}: what we believe about $\mu$ \textit{before} seeing data \\
$p(\mathbf{x} \mid \mu)$ & = & \textbf{Likelihood}: probability of the data given $\mu$ \\
$p(\mathbf{x})$ & = & \textbf{Evidence} (marginal likelihood): $\displaystyle\int_{-\infty}^{\infty} p(\mathbf{x} \mid \mu)\,p(\mu)\,d\mu$ \\
$p(\mu \mid \mathbf{x})$ & = & \textbf{Posterior}: what we believe about $\mu$ \textit{after} seeing data
\end{tabular}
\end{center}

The evidence $p(\mathbf{x})$ is a constant (it doesn't depend on $\mu$ --- it's already integrated out). So for the purpose of finding the MAP, we can ignore it:
\[
  p(\mu \mid \mathbf{x}) \propto \underbrace{p(\mathbf{x} \mid \mu)}_{\text{likelihood}} \cdot \underbrace{p(\mu)}_{\text{prior}}
\]

\textbf{In words:} posterior $\propto$ likelihood $\times$ prior. The MAP estimate is the value of $\mu$ that maximizes this product.

\subsubsection*{Step 2: Write out each piece explicitly}

\textbf{Likelihood} (from Part a):
\[
  p(\mathbf{x} \mid \mu) = \prod_{i=1}^n \frac{1}{\sqrt{2\pi}\,\sigma}\exp\!\left(-\frac{(x_i - \mu)^2}{2\sigma^2}\right)
\]

\textbf{Prior:}
\[
  p(\mu) = \frac{1}{\sqrt{2\pi}\,\tau}\exp\!\left(-\frac{(\mu - m)^2}{2\tau^2}\right)
\]

\textbf{Log-posterior} (taking $\ln$ of the product, dropping all terms that don't involve $\mu$):
\begin{align*}
  \ln p(\mu \mid \mathbf{x}) &= \ln p(\mathbf{x} \mid \mu) + \ln p(\mu) + \text{const} \\[6pt]
  &= \underbrace{-\frac{1}{2\sigma^2}\sum_{i=1}^n(x_i - \mu)^2}_{\text{log-likelihood (data term)}} \;\;\underbrace{- \frac{(\mu - m)^2}{2\tau^2}}_{\text{log-prior (regularizer)}} + \text{const}
\end{align*}

\textbf{Notice the structure:} the log-posterior is a sum of two quadratic penalties on $\mu$:
\begin{itemize}
  \item The likelihood pulls $\mu$ toward the data (toward $\bar{x}$).
  \item The prior pulls $\mu$ toward the prior mean $m$.
  \item The MAP estimate will balance these two forces.
\end{itemize}

\subsubsection*{Step 3: Differentiate and solve}

Take the derivative with respect to $\mu$ and set it to zero:
\[
  \frac{d}{d\mu}\ln p(\mu \mid \mathbf{x}) = \frac{1}{\sigma^2}\sum_{i=1}^n(x_i - \mu) - \frac{1}{\tau^2}(\mu - m) \stackrel{!}{=} 0
\]

Expand the first term using $\sum(x_i - \mu) = n\bar{x} - n\mu$:
\[
  \frac{n\bar{x} - n\mu}{\sigma^2} - \frac{\mu - m}{\tau^2} = 0
\]

Distribute:
\[
  \frac{n\bar{x}}{\sigma^2} - \frac{n\mu}{\sigma^2} - \frac{\mu}{\tau^2} + \frac{m}{\tau^2} = 0
\]

Move all $\mu$ terms to the left, everything else to the right:
\[
  \mu\cdot\frac{n}{\sigma^2} + \mu\cdot\frac{1}{\tau^2} = \frac{n\bar{x}}{\sigma^2} + \frac{m}{\tau^2}
\]

Factor out $\mu$:
\[
  \mu\!\left(\frac{n}{\sigma^2} + \frac{1}{\tau^2}\right) = \frac{n\bar{x}}{\sigma^2} + \frac{m}{\tau^2}
\]

Divide both sides:
\[
  \boxed{\hat\mu_{\text{MAP}} = \frac{\dfrac{n}{\sigma^2}\,\bar{x} \;+\; \dfrac{1}{\tau^2}\,m}{\dfrac{n}{\sigma^2} + \dfrac{1}{\tau^2}}}
\]

\subsubsection*{Step 4: Interpret --- it's a weighted average}

Rewrite the formula as:
\[
  \hat\mu_{\text{MAP}} = w_{\text{data}}\,\bar{x} + w_{\text{prior}}\,m, \qquad w_{\text{data}} + w_{\text{prior}} = 1
\]

where:
\[
  w_{\text{data}} = \frac{n/\sigma^2}{n/\sigma^2 + 1/\tau^2}, \qquad
  w_{\text{prior}} = \frac{1/\tau^2}{n/\sigma^2 + 1/\tau^2}
\]

What determines the weights? The quantities $n/\sigma^2$ and $1/\tau^2$ are \textbf{precisions} (inverse variances):

\begin{itemize}
  \item \textbf{Data precision} = $n/\sigma^2$: how precisely the data pins down $\mu$. Grows with $n$ (more readings) and shrinks with $\sigma^2$ (noisier sensor).
  \item \textbf{Prior precision} = $1/\tau^2$: how confident we are in the prior belief. A tight prior ($\tau$ small) means high precision; a vague prior ($\tau$ large) means low precision.
\end{itemize}

\textbf{The rule is simple:} whichever source of information is more precise gets more weight in the final estimate.

\subsubsection*{Step 5: Compute}

Plug in $n = 5$, $\sigma^2 = 4$, $\bar{x} = 21.34$, $m = 20$, $\tau^2 = 9$:

Data precision: $\dfrac{n}{\sigma^2} = \dfrac{5}{4} = 1.25$

\medskip
Prior precision: $\dfrac{1}{\tau^2} = \dfrac{1}{9} \approx 0.111$

\medskip
Total precision: $1.25 + 0.111 = 1.361$

\medskip
Weights:
\[
  w_{\text{data}} = \frac{1.25}{1.361} = 0.918, \qquad w_{\text{prior}} = \frac{0.111}{1.361} = 0.082
\]

\[
  \boxed{\hat\mu_{\text{MAP}} = 0.918 \times 21.34 + 0.082 \times 20 = 19.59 + 1.64 = 21.23\;{}^\circ\text{C}}
\]


% ----------------------------------------
\subsection*{Part (c): How much does the prior pull?}

The MLE gives $21.34$ and the MAP gives $21.23$, so the prior pulls the estimate by:
\[
  21.34 - 21.23 = 0.11\;{}^\circ\text{C} \quad\text{toward the prior mean of } 20
\]

This is a small shift. Why? Look at the weights: the data gets 91.8\% of the weight, the prior only 8.2\%.

\textbf{Intuition:} The data precision ($n/\sigma^2 = 1.25$) is much larger than the prior precision ($1/\tau^2 = 0.111$). In other words, the prior is quite vague ($\tau = 3$ is wide) while 5 measurements with $\sigma = 2$ already pin down $\mu$ reasonably well. The data ``speaks louder'' than the prior.

\medskip
Another way to see it: the prior says ``$\mu$ is somewhere in $(20 \pm 6)$ with 95\% probability'' ($\tau = 3$, so the 95\% interval is $m \pm 2\tau$). The data, with $\text{SE} = \sigma/\sqrt{n} = 2/\sqrt{5} \approx 0.89$, says ``$\mu$ is in $(21.34 \pm 1.75)$.'' The data's interval is much narrower, so it dominates.

\begin{center}
\renewcommand{\arraystretch}{1.4}
\begin{tabular}{lcc}
\hline
& \textbf{Mean} & \textbf{Precision (= $1/\text{Var}$)} \\
\hline
Prior & 20.00 & $1/9 = 0.111$ \\
Data (likelihood) & 21.34 & $5/4 = 1.250$ \\
\textbf{MAP (combined)} & \textbf{21.23} & $\mathbf{1.361}$ \\
\hline
\end{tabular}
\end{center}

\textbf{Key insight:} The MAP precision is always \textit{greater} than either the data or prior precision alone. Combining information always makes you more certain, never less.


% ----------------------------------------
\subsection*{Part (d): What happens with $n = 100$ readings?}

Now suppose we take $n = 100$ readings with the same sample mean $\bar{x} = 21.34$.

Data precision: $\dfrac{n}{\sigma^2} = \dfrac{100}{4} = 25$

\medskip
Prior precision: $\dfrac{1}{\tau^2} = \dfrac{1}{9} = 0.111$

\medskip
New weights:
\[
  w_{\text{data}} = \frac{25}{25.111} = 0.9956, \qquad w_{\text{prior}} = \frac{0.111}{25.111} = 0.0044
\]

Or directly:
\[
  \hat\mu_{\text{MAP}} = \frac{25 \times 21.34 + 0.111 \times 20}{25.111} = \frac{533.5 + 2.22}{25.111} = \frac{535.72}{25.111} = 21.33\;{}^\circ\text{C}
\]

Compare:
\begin{center}
\renewcommand{\arraystretch}{1.4}
\begin{tabular}{lccc}
\hline
& \textbf{MLE} & \textbf{MAP ($n=5$)} & \textbf{MAP ($n=100$)} \\
\hline
Estimate & 21.34 & 21.23 & 21.33 \\
$w_{\text{data}}$ & 100\% & 91.8\% & 99.6\% \\
$w_{\text{prior}}$ & 0\% & 8.2\% & 0.4\% \\
\hline
\end{tabular}
\end{center}

With 100 readings, the MAP is virtually identical to the MLE. The prior's influence has shrunk from 8.2\% to 0.4\%.

\bigskip
\textbf{The lesson:}
\begin{itemize}
  \item With \textit{little} data, the prior matters --- it stabilizes the estimate when the data alone is noisy.
  \item With \textit{lots} of data, the prior washes out --- MAP $\to$ MLE. The data overwhelms any prior belief.
  \item Mathematically: as $n \to \infty$, $w_{\text{data}} = \frac{n/\sigma^2}{n/\sigma^2 + 1/\tau^2} \to 1$, so $\hat\mu_{\text{MAP}} \to \bar{x} = \hat\mu_{\text{MLE}}$.
  \item This is a fundamental property of Bayesian estimation: \textbf{the posterior is asymptotically dominated by the likelihood}. No matter what prior you start with (as long as it's nonzero at the true value), enough data will override it.
\end{itemize}

\end{document}
